% Original-author contribution to planned164.4; distinct from Musgrave2011.
\chapter{Popper and Hypothetico-Deductive Testing}
\label{sec:p15-popper-hypothetico-deductive-reasoning}

Deduction preserves truth from premises to consequences. It does not
establish the empirical premises, identify a unique culprit when a
test fails, or verify a theory by affirming its predicted consequence.
Hypothetico-deductive testing needs all three distinctions.

The primary text here is Popper's \emph{The Logic of Scientific
Discovery}, Routledge Classics 2002, electronically issued in 2005.
Its translators' note identifies the author-assisted 1959 English
translation of the German book published in autumn 1934 with imprint
1935. Starred notes and dated addenda include later corrections.
This is not an inspected German first edition. In particular,
the 1972 account of rational preference is not dated to 1934
\parencite[translators' note; \S85, addendum]{Popper2002LogicDiscoveryUI}.
The philosophical arguments retain their attributed status. The annex
results prove their stated claims, not the unique correctness
of a demarcation criterion. Popper's probability and quantum chapters,
with their later notes, are treated in \cref{sec:p15-w02b-p15-027}.

\section{The induction argument and the task of enquiry}
\label{sec:p15-popper-induction}

Popper's \S1 challenges a foundation which infers universal theories
from singular experiential reports. A synthetic principle of induction
would itself need justification. Grounding it in experience employs
the contested passage again; requiring a higher principle repeats the
demand. An analytic principle cannot supply the additional empirical
restriction. Replacing truth by probability does not, without further
assumptions, justify the transfer
\parencite[\S1]{Popper2002LogicDiscoveryUI}.
This is his argument against the proposed foundation, not a theorem
that every statistical or Bayesian model is inconsistent.

The questions how someone invents a hypothesis, what follows from it,
and why it should guide action are distinct. A psychological history
does not establish a deduction; a deduction does not establish a
preference. Popper's \S\S2--3 separate invention from examination,
and examination into consistency, empirical form, comparison with
rivals and tests of applications. This does not require invention
without observation: it identifies the logical reconstruction's task.

The Hume treatment in \cref{sec:p15-hume-and-the-problem-of-induction}
distinguishes nonentailment from a further demand for a foundation.
Whether a countermodel defeats a nondeductive warrant depends on
what that warrant claims. Naming every such warrant induction is not
a refutation by itself. Goodman's treatment in
\cref{sec:p15-goodman-projectibility} likewise distinguishes the
choice of projectible predicates from the support assigned to a chosen
candidate. No mere count of favourable cases settles both questions.

\section{Theory, auxiliaries and initial conditions}
\label{sec:p15-popper-test-package}

\begin{hbdefinition}{def:p15-popper-test-record}{A conditional test record}
In classical two-valued semantics, let \(H\) represent the target
theory, \(A\) its test auxiliaries, \(I\) the initial-condition
description, and \(E\) the predicted observational statement.
Auxiliaries may concern calibration, specimen identification,
detector response and interference; conditions concern this application.
The record specifies \(H\land A\land I\models E\), empirical grounds
for its components, the outcome description, comparison theories and date.

The sign \(\models\) denotes semantic consequence.
The formula \((H\land A\land I)\to E\) is an object-language
conditional, not the assertion that it is valid. Popper's later
starred correction to \S18 acknowledges this distinction.
An accepted failed-outcome statement is represented here as \(\neg E\);
perception does not thereby become an infallible logical negation.
\end{hbdefinition}

Popper's \S12 deduces predictions from laws and conditions; \S18
explicitly assigns the contradiction to the entire premise system.
Combining theories and conditions into one system does not select
one conjunct as the unique failure
\parencite[\S\S12,18]{Popper2002LogicDiscoveryUI}.

\begin{hbtheorem}{thm:p15-popper-package}{proved=annex:p15.popper.package}{Package refutation and premise isolation}
For independently assignable propositional atoms put
\(\chi=((H\land A\land I)\to E)\). Then
\[
 \chi,\neg E\models\neg(H\land A\land I),\qquad
 \chi,A,I,\neg E\models\neg H.
\]
The first two premises do not entail \(\neg H,\neg A,\neg I\)
separately. Nor do \(\chi,A,I,E\) entail \(H\).
The entailments persist under substitution of arbitrary classical
formulas; the nonentailments concern the indicated independent atoms.
\end{hbtheorem}

The second deduction holds \(A,I\) fixed as true. Their empirical
acceptance has grounds not supplied by a truth table.
Retesting an instrument or a condition estimate is therefore
compatible with modus tollens. Merely mentioning possible error,
however, does not explain an anomaly.
The distinction refines the common package boundary in
\cref{thm:p11-rat141-test-logic}.

\begin{hbexample}{def:p15-popper-thread}{A complete idealized test}
On a nonempty specimen domain \(U\), positive real-valued functions
\(L,S\) denote actual load and breaking threshold; \(r_L,r_S\)
are their instrument reports in one fixed unit.
Numerals and the order \(>\) have their usual real interpretations.
The predicates \(\mathsf{Br},\mathsf R\) denote actual and reported
breaking. This is an idealized deterministic model, not a claim
about all real materials. Let
\[
\begin{aligned}
H&:\ \forall x(L(x)>S(x)\to\mathsf{Br}(x)),\\
A&:\ \forall x\bigl(r_L(x)=L(x)\land r_S(x)=S(x)
             \land(\mathsf R(x)\leftrightarrow\mathsf{Br}(x))\bigr),\\
I&:\ r_L(c)=2\land r_S(c)=1,\qquad E:\mathsf R(c).
\end{aligned}
\]
Calibration turns recorded readings into physical conditions.
The law then gives actual breaking and the detector gives \(E\).
Suppose the accepted outcome instead is \(\neg\mathsf R(c)\).

The numbers follow the structure of Popper's thread illustration.
His later \S12 note distinguishes a general law, a structure-specific
threshold law and two conditions. Reporting functions and the
following models are handbook additions
\parencite[\S12,note *1]{Popper2002LogicDiscoveryUI}.
\end{hbexample}

\begin{hbtheorem}{thm:p15-popper-thread}{proved=annex:p15.popper.thread}{The same failure report admits different diagnoses}
\(H,A,I\models E\). Three one-specimen interpretations satisfy
\(I,\neg E\): one with \(A\) true and \(H\) false; one with
\(H\) true and false strength calibration; and one with \(H\)
true and false detector equivalence.
The complete displayed observation record does not identify the
failed component.
\end{hbtheorem}

\begin{center}\small
\begin{tabular}{@{}lccccc@{}}
\toprule
Diagnosis&\(L(c)\)&\(S(c)\)&\(\mathsf{Br}(c)\)&\(\mathsf R(c)\)&Failure\\
\midrule
Law&2&1&false&false&\(H\)\\
Calibration&2&3&false&false&strength clause of \(A\)\\
Detector&2&1&true&false&detector clause of \(A\)\\
\bottomrule
\end{tabular}
\end{center}

Every row has \(r_L(c)=2,r_S(c)=1\).
Independent calibration can discriminate the second row;
another checked detector or image can discriminate the third.
Those additional records themselves require testing and interpretation.
A possible countermodel is not automatically the correct account
of a laboratory.

\section{Universal, existential and bounded claims}
\label{sec:p15-popper-asymmetry}

The strict law in Popper's asymmetry ranges beyond enumerated cases.
Numerical universality instead concerns an exhaustively bounded
collection. An isolated strict existential claim makes an unbounded
existence assertion, while an existential claim in a richer theory
can add testable consequences. A witness and an unlimited unsuccessful
search have different force \parencite[\S\S13--15, including
the starred qualification]{Popper2002LogicDiscoveryUI}.

\begin{hbtheorem}{thm:p15-popper-finite-record}{proved=annex:p15.popper.finite-record}{Finite records and exhaustive-domain premises}
Let \(n\geq0\), with distinct named objects \(a_1,\ldots,a_n\)
and unary predicates \(\mathsf{Rv},\mathsf{Wh}\).
Let \(D_n\) assert \(\mathsf{Rv}(a_i)\land\mathsf{Wh}(a_i)\)
for every \(i\), together with pairwise distinctness, and put
\(H_{\rm rw}=\forall x(\mathsf{Rv}(x)\to\mathsf{Wh}(x))\).
Two nonempty finite structures agree on all indicated predicate
facts about the named objects and satisfy \(D_n\), but one satisfies
\(H_{\rm rw}\) and the other its negation.
Nevertheless
\(\mathsf{Rv}(b)\land\neg\mathsf{Wh}(b)\models\neg H_{\rm rw}\).

For another unary predicate \(P\),
\(\bigwedge_{i=1}^n\neg P(a_i)\not\models\neg\exists xP(x)\),
whereas \(P(b)\models\exists xP(x)\).
For \(n\geq1\), relative to
\(\forall x\bigvee_{i=1}^n x=a_i\), \(H_{\rm rw}\) is equivalent to
\(\bigwedge_{i=1}^n(\mathsf{Rv}(a_i)\to\mathsf{Wh}(a_i))\).
\end{hbtheorem}

Treating a sample as an enumeration silently adds the domain premise.
Applying species, colour or location predicates is an empirical
identification issue, not another quantifier rule. A mistaken report
can leave a valid conditional refutation inapplicable.

\section{Demarcation, content and evasion}
\label{sec:p15-popper-falsifiers}

Possible basic statements are not a jointly accepted record:
their class includes incompatible possibilities. Falsifiability
asks what a theory prohibits; empirical falsification additionally
concerns accepted refuting statements and a tested falsifying
hypothesis. Demarcation is a methodological proposal, not a criterion
of meaning or a theorem assigning value to every nonscientific claim
\parencite[\S\S4,6,21--22]{Popper2002LogicDiscoveryUI}.

\begin{hbdefinition}{def:p15-popper-relative-falsifiers}{A declared observational comparison}
Fix a nonempty class \(\mathcal K\) of classical interpretations
of one fixed language
and a nonempty set \(\mathcal B\) of possible test statements,
each satisfiable in \(\mathcal K\). They declare the shared
background and observational vocabulary, not necessarily every
scientific observation. Define
\[
F_{\mathcal K,\mathcal B}(T)
=\{b\in\mathcal B:\text{no }\mathcal M\in\mathcal K
                         \text{ satisfies }T\land b\}.
\]
For a nondegenerate comparison require both this class and its
complement in \(\mathcal B\) nonempty.
Acceptance of an actual record is not part of the definition.
\end{hbdefinition}

\begin{hbtheorem}{thm:p15-popper-falsifier-order}{proved=annex:p15.popper.falsifier-order}{A partial ordering of forbidden tests}
If every \(\mathcal K\)-model of \(T_1\) satisfies \(T_2\),
then \(F_{\mathcal K,\mathcal B}(T_2)\subseteq
F_{\mathcal K,\mathcal B}(T_1)\).
Inclusion is strict if some admitted test is compatible with
\(T_2\) in \(\mathcal K\) but incompatible with \(T_1\).
Strict logical strengthening alone need not give strict inclusion
in an incomplete observational vocabulary.
A tautology has no potential falsifiers and a contradiction all
of \(\mathcal B\). Falsifier classes can be incomparable.
\end{hbtheorem}

This modern relative illustration of \S\S31--35 supplies no
universal numerical severity scale. Inclusion orders some classes,
not all; infinite cardinalities need not capture the desired
comparison. Changing background or outcome vocabulary changes
what the comparison establishes.
Replacing a precise prediction by its disjunction with its contrary
removes that prediction's possible refutation. Whether to accept
that weakening is a different question.

Popper grants the logical possibility of saving a theory by changing
auxiliaries, definitions or accepted observations; his reply constrains
how science should treat revisions. It does not make revision
mechanically decidable. The later starred \S20 note withdraws the
assertion that the Lorentz--Fitzgerald contraction hypothesis had
no falsifiable consequences. That example cannot be repeated as
established total untestability
\parencite[\S\S19--20,note *1]{Popper2002LogicDiscoveryUI}.
Nor must every auxiliary correction be rejected: a miscalibration
can be independently investigated. Record the changed system,
derive further consequences and specify its countertests.

\section{The empirical basis and a stopping decision}
\label{sec:p15-popper-basic-statements}

A basic statement in \S28 reports a testable occurrence in an
individual space--time region, not merely a conviction.
Testing can continue through further statements, but completing
an investigation requires provisional acceptance. Experience
motivates that decision; logical justification is another relation.
Agreement does not make the statement true
\parencite[\S\S25--30]{Popper2002LogicDiscoveryUI}.
For the thread record, another image analysis, calibration or
clock check can reopen the accepted description.

A demand for final proof of every empirical premise cannot stop
without further premises. Revisable testing can stop for a purpose
without declaring the stopping statement immune to criticism.
This does not by itself justify every practical reliance.
Reproducibility must not require literal repetition of every unique
past occurrence: Popper's later \S22 note appeals also to publicly
testable witnesses and records.

Hermann's 1935 review grants revisable theories and theory-laden
observations but objects that moving unjustified decisions between
statements does not explain science's claim to knowledge.
The available text is Bacciagaluppi's 2020 English translation,
not an inspected German printing. Popper's \S\S29--30 supply
an opposing position to assess, not evidence of a direct reply
to that review \parencite[pp.~621--623]{Hermann2020PopperReview}.

Public correction can constrain decisions without guaranteeing truth.
The disagreement concerns whether epistemic authority needs an
ultimate justification, not whether observations are either
infallible or useless.

\section{Corroboration is not truth or posterior probability}
\label{sec:p15-popper-corroboration}

Corroboration appraises an identified history of tests and criticism,
not mere compatibility. It attends to severe attempted refutations
and an accepted record up to a stated time. Popper's 1972 addendum
distinguishes risky rational preference from justified belief in
truth or probable truth, and denies that corroboration measures
verisimilitude \parencite[\S\S82--85]{Popper2002LogicDiscoveryUI}.
Easy predictions shared by every rival do not discriminate them.
An independently calibrated test on which rivals prescribe
incompatible outcomes can discriminate, without proving the survivor.

\begin{hbtheorem}{thm:p15-popper-probability-order}{proved=annex:p15.popper.probability-order}{Nested hypotheses under a specified probability law}
In a probability space \((\Omega,\mathcal F,\mathbb P)\), let measurable
\(V_1\subseteq V_2\), and let \(D\) be measurable with
\(\mathbb P(D)>0\). Then
\[
\begin{aligned}
\mathbb P(V_1\mid D)&\leq\mathbb P(V_2\mid D),\\
\mathbb P(V_2\mid D)-\mathbb P(V_1\mid D)
&=\frac{\mathbb P((V_2\setminus V_1)\cap D)}{\mathbb P(D)}.
\end{aligned}
\]
Strict set inclusion need not give a strict conditional inequality.
\end{hbtheorem}

This modern illustration identifies a tension behind \S83:
preferring an informative restriction is not maximizing the probability
of a containing event. It does not make Bayesian comparison of
mutually exclusive models impossible, give every universal law zero
probability under every prior, or identify Popper's different
numerical proposals. The treatments beginning at
\cref{sec:p15-bayesian-updating-prediction} specify their own model
and prior; their quantities are not renamed corroboration here.

\begin{hbtheorem}{thm:p15-popper-coin}{proved=annex:p15.popper.coin}{A rare outcome is not a logical counterexample to a chance law}
For \(n\geq1\), \(\Omega=\{0,1\}^n\) with uniform mass assigns
each sequence probability \(2^{-n}>0\). The all-one sequence is
possible; at \(n=10\) its probability is \(1/1024<0.05\).
Its occurrence is consistent with that law even if a declared
decision procedure rejects on the outcome.
\end{hbtheorem}

Low probability differs from prohibition. Statistical rejection
needs a sampling model and error criterion; surprise does not make
it deductive falsification. Possible realization under the law also
does not prove that the laboratory mechanism satisfies the law.

\section{The operative method and its philosophical limit}
\label{sec:p15-popper-method}

Identify target and rivals; record auxiliaries and conditions;
derive competing outcomes; arrange checks sensitive to their
difference; specify scrutiny of observations and apparatus.
Report both deduction and accepted empirical premises.
After failure compare law, auxiliary and condition diagnoses.
After survival report tests survived and alternatives discriminated.
A revision needs a new record, not retrospective alteration
of the prediction.

The auxiliary objection is logically correct; the proposed reply
exposes the revised system to tests rather than selecting blame
by logic alone. The existential objection requires distinguishing
bounded search and contextual consequences from unlimited existence.
The objection that achievements are ignored is answered by comparative
demanding tests, not affirming the consequent.

The remaining induction question asks why past performance warrants
trust in a new application. Changing the goal to critical selection
does not automatically provide a justified future-success probability.
Comparison with confirmation and inductive-probability accounts
therefore concerns different proposed warrants, not a deductive
victory obtained by defining their conclusions as non-deductions.

\begin{proofboundary}{def:p15-popper-source-boundary}{The exact scope of this treatment}
The formal results establish only their declared classical, finite
or probabilistic claims. They do not prove the philosophical rule of
maximizing criticism or resolve every objection. In the probability
chapter, the third binomial formula and the derivation of Bernoulli's
theorem from the two requirements of \S64 remain Popper's sketch and
claim (\cref{sec:p15-w02b-p15-027}); the constructions of
appendix~iv, the appendices~v--vii on the quantum examples, the new
appendices~*ii--*xii and the \emph{Postscript} are not reconstructed.
The quantum chapter contains no derivation of the Heisenberg formulae,
and none is given. Later notes and 1972 addenda retain their dates
(\cref{sec:p15-w02b-p15-027-later-notes}).
Musgrave's separate complete 2011 \emph{Popper and
Hypothetico-deductivism}, \emph{Handbook of the History of Logic},
volume~10, pp.~205--234, is not replaced by this original-author
treatment or by Hermann's review.
\end{proofboundary}

\index{Popper, Karl}\index{hypothetico-deductive method}
\index{falsification!auxiliary assumptions}\index{demarcation}
\index{corroboration}\index{basic statements}
